\chapter*{Processor Dependent Code}
\addcontentsline{toc}{chapter}{Processor Dependent Code}
This chapter presents the Twin-64 processor dependent code....

\section*{Concepts}
\addcontentsline{toc}{section}{Concepts}
´
\section*{PDC Address Space}
\addcontentsline{toc}{section}{PDC Address Space}

0xF000\_0000 - start of PDC. Reset vector. 

0xF000\_0004 - PDC entry point. Jumps to PDC entry function.

0xF000\_0008 - reserved.

0xF000\_000C - reserved.

Calling convention:

ARG0 - entry point number.

ARG1, ARg2, ARG3 - arguments.

All also function as RET0 to RET3.

REGn - REGm - caller saves are used locally.

PDC has perhaps a memory area in page zero to use as working memory.

PDC also needs a stack to run on -> in lower memory.

PDC is called privileged, interrupts disabled, uncached.


 

\section*{Design Rationale}
\addcontentsline{toc}{section}{Design Rationale}

\chapter*{Processor Dependent Code Functions}
\addcontentsline{toc}{chapter}{Processor Dependent Code}




